% ==============================================================================
% LIBRERÍA DE DIAGRAMAS TIKZ PARA CASOS DE USO (SIN MARCO) - TT 2026-B009
% ==============================================================================

% --- ESTILO BASE PARA LOS CASOS DE USO ---
\tikzset{
    actor/.style={thick},
    sistema/.style={thick, fill=gray!5},
    caso/.style={draw, ellipse, fill=azulESCOM!15, minimum width=3.5cm, minimum height=0.8cm, align=center},
    relacion/.style={thick},
    extension/.style={thick, dashed, <-, >=stealth}
}

% --- COMANDO INTERNO PARA DIBUJAR AL ACTOR (Stickman) ---
\newcommand{\dibujarQuejoso}[2]{ % #1: x, #2: y
    \draw[actor] (#1,#2+0.7) circle (0.2); 
    \draw[actor] (#1,#2+0.5) -- (#1,#2); 
    \draw[actor] (#1-0.3,0.3+#2) -- (#1+0.3,0.3+#2); 
    \draw[actor] (#1,#2) -- (#1-0.2,#2-0.3); 
    \draw[actor] (#1,#2) -- (#1+0.2,#2-0.3); 
    \node at (#1,#2-0.6) {Quejoso};
}

% --- DEFINICIÓN DE LAS 18 FIGURAS (SIN MARCO) ---

% MODULO PERFIL
\newcommand{\figCUuno}{ \begin{tikzpicture}[font=\small] \dibujarQuejoso{-3.2}{2} \draw[sistema] (-1,0.5) rectangle (5,3.5); \node[caso] (c) at (2,2) {CU01: Crear\\cuenta}; \draw[relacion] (-2.9,2.3) -- (c); \end{tikzpicture} }

\newcommand{\figCUdos}{ \begin{tikzpicture}[font=\small] \dibujarQuejoso{-3.2}{2} \draw[sistema] (-1,0.5) rectangle (5,3.5); \node[caso] (c) at (2,2) {CU02: Iniciar\\sesión}; \draw[relacion] (-2.9,2.3) -- (c); \end{tikzpicture} }

\newcommand{\figCUtres}{ \begin{tikzpicture}[font=\small] \dibujarQuejoso{-3.2}{2.5} \draw[sistema] (-1,0) rectangle (5,4); \node[caso] (c2) at (2,3) {CU02: Iniciar sesión}; \node[caso] (c3) at (2,1) {CU03: Recuperar clave}; \draw[extension] (c2) -- (c3) node[midway, right] {\scriptsize $<<Extend>>$}; \draw[relacion] (-2.9,2.8) -- (c2); \end{tikzpicture} }

\newcommand{\figCUcuatro}{ \begin{tikzpicture}[font=\small] \dibujarQuejoso{-3.2}{2} \draw[sistema] (-1,0.5) rectangle (5,3.5); \node[caso] (c) at (2,2) {CU04: Editar\\perfil}; \draw[relacion] (-2.9,2.3) -- (c); \end{tikzpicture} }

\newcommand{\figCUcinco}{ \begin{tikzpicture}[font=\small] \dibujarQuejoso{-3.2}{2} \draw[sistema] (-1,0.5) rectangle (5,3.5); \node[caso] (c) at (2,2) {CU05: Registrar\\tutor}; \draw[relacion] (-2.9,2.3) -- (c); \end{tikzpicture} }

% MODULO QUEJAS
\newcommand{\figCUseis}{ \begin{tikzpicture}[font=\small] \dibujarQuejoso{-3.2}{2} \draw[sistema] (-1,0.5) rectangle (5,3.5); \node[caso] (c) at (2,2) {CU06: Levantar\\queja}; \draw[relacion] (-2.9,2.3) -- (c); \end{tikzpicture} }

\newcommand{\figCUsiete}{ \begin{tikzpicture}[font=\small] \dibujarQuejoso{-3.2}{2} \draw[sistema] (-1,0.5) rectangle (5,3.5); \node[caso] (c) at (2,2) {CU07: Adjuntar\\evidencia}; \draw[relacion] (-2.9,2.3) -- (c); \end{tikzpicture} }

\newcommand{\figCUocho}{ \begin{tikzpicture}[font=\small] \dibujarQuejoso{-3.2}{2} \draw[sistema] (-1,0.5) rectangle (5,3.5); \node[caso] (c) at (2,2) {CU08: Consultar\\tablero}; \draw[relacion] (-2.9,2.3) -- (c); \end{tikzpicture} }

\newcommand{\figCUnueve}{ \begin{tikzpicture}[font=\small] \dibujarQuejoso{-3.2}{2} \node[caso] (c6) at (1.2,2) {CU06: Levantar\\queja}; \node[caso] (c9) at (5.2,2) {CU09: Modificar\\queja}; \draw[extension] (c6) -- (c9) node[midway, above] {\scriptsize $<<Extend>>$}; \draw[relacion] (-2.9,2.3) -- (c6); \end{tikzpicture} }

\newcommand{\figCUdiez}{ \begin{tikzpicture}[font=\small] \dibujarQuejoso{-3.2}{2} \node[caso] (c6) at (1.2,2) {CU06: Levantar\\queja}; \node[caso] (c10) at (5.2,2) {CU10: Eliminar\\borrador}; \draw[extension] (c6) -- (c10) node[midway, above] {\scriptsize $<<Extend>>$}; \draw[relacion] (-2.9,2.3) -- (c6); \end{tikzpicture} }

\newcommand{\figCUonce}{ \begin{tikzpicture}[font=\small] \dibujarQuejoso{-3.2}{2} \draw[sistema] (-1,0.5) rectangle (5,3.5); \node[caso] (c) at (2,2) {CU11: Recibir\\notificaciones}; \draw[relacion] (-2.9,2.3) -- (c); \end{tikzpicture} }

% MODULO SEGUIMIENTO
\newcommand{\figCUdoce}{ \begin{tikzpicture}[font=\small] \dibujarQuejoso{-3.2}{2} \node[caso] (c11) at (1.2,2) {CU11: Recibir\\notif.}; \node[caso] (c12) at (5.2,2) {CU12: Completar\\info.}; \draw[extension] (c11) -- (c12) node[midway, above] {\scriptsize $<<Extend>>$}; \draw[relacion] (-2.9,2.3) -- (c11); \end{tikzpicture} }

\newcommand{\figCUtrece}{ \begin{tikzpicture}[font=\small] \dibujarQuejoso{-3.2}{2} \node[caso] (c11) at (1.2,2) {CU11: Recibir\\notif.}; \node[caso] (c13) at (5.2,2) {CU13: Consultar\\medidas}; \draw[extension] (c11) -- (c13) node[midway, above] {\scriptsize $<<Extend>>$}; \draw[relacion] (-2.9,2.3) -- (c11); \end{tikzpicture} }

\newcommand{\figCUcatorce}{ \begin{tikzpicture}[font=\small] \dibujarQuejoso{-3.2}{2} \node[caso] (c11) at (1.2,2) {CU11: Recibir\\notif.}; \node[caso] (c14) at (5.2,2) {CU14: Propuesta\\conciliación}; \draw[extension] (c11) -- (c14) node[midway, above] {\scriptsize $<<Extend>>$}; \draw[relacion] (-2.9,2.3) -- (c11); \end{tikzpicture} }

\newcommand{\figCUquince}{ \begin{tikzpicture}[font=\small] \dibujarQuejoso{-3.2}{2} \node[caso] (c14) at (1.2,2) {CU14: Consultar\\propuesta}; \node[caso] (c15) at (5.2,2) {CU15: Aceptar\\propuesta}; \draw[extension] (c14) -- (c15) node[midway, above] {\scriptsize $<<Extend>>$}; \draw[relacion] (-2.9,2.3) -- (c14); \end{tikzpicture} }

\newcommand{\figCUdieciseis}{ \begin{tikzpicture}[font=\small] \dibujarQuejoso{-3.2}{2} \node[caso] (c14) at (1.2,2) {CU14: Consultar\\propuesta}; \node[caso] (c16) at (5.2,2) {CU16: Rechazar\\propuesta}; \draw[extension] (c14) -- (c16) node[midway, above] {\scriptsize $<<Extend>>$}; \draw[relacion] (-2.9,2.3) -- (c14); \end{tikzpicture} }

\newcommand{\figCUdiecisiete}{ \begin{tikzpicture}[font=\small] \dibujarQuejoso{-3.2}{2} \draw[sistema] (-1,0.5) rectangle (5,3.5); \node[caso] (c) at (2,2) {CU17: Descargar\\acuse}; \draw[relacion] (-2.9,2.3) -- (c); \end{tikzpicture} }

\newcommand{\figCUdieciocho}{ \begin{tikzpicture}[font=\small] \dibujarQuejoso{-3.2}{2} \draw[sistema] (-1,0.5) rectangle (5,3.5); \node[caso] (c) at (2,2) {CU18: Detalles\\completos}; \draw[relacion] (-2.9,2.3) -- (c); \end{tikzpicture} }